\chapter{贝叶斯机器学习}
\begin{quote}\small
\textit{本章摘要。}本章将点预测扩展为对未知量及其依赖关系的分布表达。先建立贝叶斯更新和后验预测，再用概率图模型描述设备隐状态、观测和相邻标签之间的关系。Hidden Markov Model (HMM)、最大熵马尔科夫模型、Conditional Random Field (CRF) 和马尔科夫随机场分别揭示生成建模、条件建模与局部依赖的不同取舍。在此基础上讨论参数后验的近似计算，最后以校准、敏感性分析和维护决策检验概率是否可信。
\end{quote}
前一部分已经固定数据可用时间、预测目标与验证规则。现在假设这些前提都成立，仍然存在一个未解决的问题：少量故障能支持多精确的风险估计，多个解释相近的参数应如何共同参与预测？本章不改变任务，而是把单一答案扩展成对可能答案及其依赖关系的分布。

阅读时可以把未知量分成三类：设备长期的故障倾向属于参数，轴承此刻是否磨损属于隐状态，下一次传感器读数还包含新的随机波动。对状态求和和对参数积分处理的是不同的未知量，决定是否停机还要另外比较动作成本。先从能手算的共轭模型开始，再看链式图模型怎样减少枚举，最后讨论近似计算的误差。使用贝叶斯方法以后，仍要检查假设是否合适、概率是否校准，不能省去验证。

\section{贝叶斯更新与后验预测}

% auxiliary-resource-guide
本章末的“辅助学习材料”列出与核心方法对应的讲解和实践入口，可在阅读相关推导后，按所列小节对照学习。

\paragraph{从参数点到参数分布。}
贝叶斯建模把未知参数视为随机变量，是概率机器学习的基本表述之一\cite{murphy2022pml,bishop2006pattern}。给定先验$p(\bm{\theta})$和似然$p(\mathcal{D}\mid\bm{\theta})$，后验为
\begin{equation}
 p(\bm{\theta}\mid\mathcal{D})=\frac{p(\mathcal{D}\mid\bm{\theta})p(\bm{\theta})}{p(\mathcal{D})},
 \qquad p(\mathcal{D})=\int p(\mathcal{D}\mid\bm{\theta})p(\bm{\theta})\,d\bm{\theta}.
\end{equation}
先验表达观测数据之前的结构知识，似然描述数据如何由参数生成，后验则是二者结合后的知识状态。

\paragraph{后验预测。}
\begin{figure}[htbp]
\centering
\begin{tikzpicture}[node distance=0.8cm and 0.75cm]
 \node[idi input, minimum width=2.1cm] (prior) {先验知识\\$p(\bm\theta)$};
 \node[idi input, right=of prior, minimum width=2.1cm] (data) {观测数据\\$\mathcal D$};
 \node[idi process, below=0.85cm of $(prior)!0.5!(data)$, minimum width=2.6cm] (like) {似然\\$p(\mathcal D\mid\bm\theta)$};
 \node[idi state, below=0.85cm of like, minimum width=2.6cm] (post) {后验\\$p(\bm\theta\mid\mathcal D)$};
 \node[idi output, below=0.85cm of post, minimum width=2.6cm] (pred) {后验预测\\$p(y_*\mid x_*,\mathcal D)$};
 \draw[idi arrow] (prior)--(like); \draw[idi arrow] (data)--(like);
 \draw[idi arrow] (like)--(post); \draw[idi arrow] (post)--(pred);
\end{tikzpicture}
\caption{贝叶斯更新把先验与观测结合为后验，再对参数不确定性积分得到预测分布。}
\label{fig:bayes-predictive-flow}
\end{figure}
图\ref{fig:bayes-predictive-flow}区分后验更新与预测两个环节：先用数据修正参数分布，再以整个后验而非单个参数点形成新观测的预测。
对新输入$x_*$，预测分布需要对参数不确定性积分：
\begin{equation}
 p(y_*\mid x_*,\mathcal{D})=\int p(y_*\mid x_*,\bm{\theta})p(\bm{\theta}\mid\mathcal{D})\,d\bm{\theta}.\label{eq:bayes-predictive-integral}
\end{equation}
式\eqref{eq:bayes-predictive-integral}中积分涉及的不确定性包含数据噪声和参数不确定性。前者称为偶然不确定性（aleatoric uncertainty），后者称为认知不确定性（epistemic uncertainty）。当设备出现训练数据未覆盖的工况时，后者尤其值得关注。

\paragraph{更新与预测各用了哪条概率规则。}
联合分布可分别写成$p(\theta,\mathcal D)=p(\mathcal D\mid\theta)p(\theta)$和$p(\theta\mid\mathcal D)p(\mathcal D)$；两式相等并除以正且有限的$p(\mathcal D)$即得Bayes公式。连续观测使用的是密度，不要求单个精确观测有正概率。若数据分批到达且给定$\theta$后条件独立，则
\begin{equation}
 p(\theta\mid\mathcal D_1,\mathcal D_2)
 \propto p(\mathcal D_2\mid\theta)p(\theta\mid\mathcal D_1).
\end{equation}
同一条记录只能更新一次；相关窗口若被误当独立批次，会使后验虚假收窄。

后验预测来自全概率公式，再使用$Y_*\perp\mathcal D\mid(x_*,\theta)$。固定设计回归把输入视作已知条件；如果$x_*$的生成机制本身依赖待推断参数，则还应对它的证据更新，不能无条件沿用$p(\theta\mid\mathcal D)$。这一细节说明概率公式总是在特定条件结构下成立。

\section{概率图模型：从依赖关系到分布分解}
单个故障率只能回答总体上有多大风险，却不能解释设备如何从正常运行进入退化状态。振动和温度往往共同受某个不可直接观测的状态影响，相邻时刻的状态又存在持续性。概率图模型以节点表示随机变量，以边和局部因子表达依赖结构，将一个庞大的联合分布拆成可以解释和计算的部分\cite{bishop2006pattern,murphy2022pml}。

有向无环图对应贝叶斯网络。若变量集合为$X_1,\ldots,X_m$，节点$X_j$的父节点集合记为$\operatorname{pa}(j)$，则
\begin{equation}
 p(X_1,\ldots,X_m)=\prod_{j=1}^{m}p(X_j\mid X_{\operatorname{pa}(j)}).
\end{equation}
没有父节点的变量使用边缘分布。这个分解规定了条件独立假设，而不是仅用箭头描绘相关性。例如，在$X\leftarrow Z\rightarrow Y$中，给定共同原因$Z$后，模型假定$X$与$Y$独立。若$Z$表示设备工况，该假设意味着知道工况后，两个传感器不再提供额外的相互依赖信息。真实设备仍可能存在共享噪声，此时需要修改观测模型。

条件化也可能引入依赖。在碰撞结构$X\rightarrow Z\leftarrow Y$中，路径原本被$Z$阻断，而观测$Z$或其后代可能使$X$与$Y$相关。因而不能简单地认为“给定更多变量总会减少依赖”。一般有向图中的这类关系由 d-separation 判断。图中的方向也不自动具有因果含义，因果解释还需要额外的领域假设或干预证据。

无向图适合表达对称的相容关系，例如相邻表面区域倾向于具有相近状态。因子图则把变量节点与因子节点分开，统一表示$p(X)\propto\prod_a\psi_a(X_a)$。它强调计算所依赖的局部分解，可以表示有向模型，也可以表示无向模型。树上的求和积消息传递能够精确计算边缘概率，最大积用于求最优联合配置；含环图的复杂度则取决于连接结构，不能只按节点数判断。

这里需要分清两类未知量。设备在某一时刻的隐状态属于随机变量推断，转移矩阵和观测分布的参数属于参数学习。概率图模型可以用最大似然估计参数，也可以为参数指定先验并计算后验。因此，“贝叶斯网络”描述图的分解形式，“贝叶斯学习”描述对参数不确定性的处理方式，两者不是同义词。

\paragraph{分解与消息的代数来源。}
按有向无环图的拓扑顺序排列变量，概率链式法则先给出$\prod_jp(X_j\mid X_1,\ldots,X_{j-1})$；只有加入“给定父节点后独立于其余先前变量”的假设，才能删去多余条件得到图分解。边因此代表被允许的依赖，不是从链式法则自动推出来的物理事实。

对树形因子图，边缘化要计算$\sum_{X\setminus X_i}\prod_a\psi_a(X_a)$。利用乘法对加法的分配律，把树的一个分支先求和，就得到消息
\begin{align}
 m_{i\to a}(x_i)&=\prod_{b\in N(i)\setminus a}m_{b\to i}(x_i),\\
 m_{a\to i}(x_i)&=\sum_{x_{a\setminus i}}\psi_a(x_a)
                  \prod_{j\in N(a)\setminus i}m_{j\to a}(x_j).
\end{align}
从叶子向内再向外传递后，$p(x_i)\propto\prod_{a\in N(i)}m_{a\to i}(x_i)$。有限离散状态下这些都是有限和；连续变量则需积分或近似。树上各分支只相交于分隔变量，所以每条证据恰好计算一次；含环时重复传递不再对应这样的独立分支消元。把求和换成最大值并保存取值索引可求最优配置，但不能由最大消息恢复边缘概率。

\section{隐马尔科夫模型与状态推断}
HMM 的前向后向、最优路径与参数估计可参照 Rabiner 的经典教程\cite{rabiner1989hmm}。下面将这些运算分别对应到在线状态识别、离线诊断和历史数据训练，避免混淆可用信息。
设$x_{1:T}$为一段设备观测，$z_t\in\{1,\ldots,K\}$为时刻$t$的隐状态。隐马尔科夫模型（HMM）假定状态只依赖上一时刻，且观测在给定当前状态后不再依赖其他状态和观测：
\begin{align}
 p(z_t\mid z_{1:t-1})&=p(z_t\mid z_{t-1}),\\
 p(x_t\mid z_{1:T},x_{1:t-1})&=p(x_t\mid z_t).
\end{align}
记初始概率为$\pi_k$，转移概率为$A_{jk}=p(z_t=k\mid z_{t-1}=j)$，观测密度或概率为$b_k(x)=p(x\mid z_t=k)$。联合分布分解为
\begin{equation}
 p(x_{1:T},z_{1:T})=\pi_{z_1}b_{z_1}(x_1)
 \prod_{t=2}^{T}A_{z_{t-1},z_t}b_{z_t}(x_t).\label{eq:bayes-hmm-joint}
\end{equation}
\begin{figure}[htbp]
\centering
\begin{tikzpicture}[>=Latex,
 state/.style={circle,draw,minimum size=10mm,fill=white},
 obs/.style={circle,draw,minimum size=10mm,fill=blue!12},
 every node/.style={font=\small}]
 \foreach \t in {1,2,3,4}{
  \node[state] (z\t) at (2.4*\t,0) {$z_{\t}$};
  \node[obs] (x\t) at (2.4*\t,-1.6) {$x_{\t}$};
  \draw[->,thick] (z\t) -- (x\t);
 }
 \foreach \a/\b in {1/2,2/3,3/4}{\draw[->,thick] (z\a) -- (z\b);}
 \node[above=3mm of z1] {$\pi$};
 \node at (6,0.65) {状态转移：$p(z_t\mid z_{t-1})$};
 \node at (6,-2.45) {观测生成：$p(x_t\mid z_t)$};
\end{tikzpicture}
\caption{HMM 的有向图结构（以四个时刻为例）。白色圆为隐状态，蓝色圆为已观测变量；横向箭头表示状态转移，向下箭头表示观测生成。箭头表示概率分解关系，不自动代表因果关系。}
\label{fig:bayes-hmm-structure}
\end{figure}
图\ref{fig:bayes-hmm-structure}可以直接读成式\eqref{eq:bayes-hmm-joint}的联合分布：每个节点贡献一个给定父节点的条件分布，首个状态贡献初始分布。给定$z_t$后，$x_t$与其他时刻的变量条件独立；但只给定观测而未固定隐状态时，各时刻观测通常仍然相关。图中的向下箭头是模型定义的生成方向，推断时则可以从观测反过来更新隐状态的概率。

每行转移概率之和为一。多通道观测可以用联合高斯分布或混合分布描述，HMM 并不要求同一时刻的各传感器通道独立。真正受限制的是跨时刻的依赖：一旦给定状态序列，观测之间便被假定条件独立。

以“正常、退化、故障”三个状态为例，转移矩阵可以表达状态持续性。禁止某条转移必须有物理依据，例如只有在明确排除维修过程的运行片段内，才可能限制故障状态不能直接返回正常。无监督学习得到的三个隐状态也不会自动具有这三个语义，需要通过标注、特征分布或专家检查建立对应关系。

\paragraph{过滤、平滑与最优路径。}
在线维护需要$p(z_t\mid x_{1:t})$，称为过滤；事后分析使用$p(z_t\mid x_{1:T})$，称为平滑。后者利用未来观测，不能直接作为实时告警的性能。若需要解释整段运行过程，还可以寻找联合概率最大的状态路径。这三个问题的输出不同，即使使用同一个 HMM，也不能混用结果。

直接枚举所有状态路径需要考虑$K^T$种配置。前向算法利用链结构复用公共计算，令$\alpha_t(k)=p(x_{1:t},z_t=k)$，则
\begin{align}
 \alpha_1(k)&=\pi_kb_k(x_1),\\
 \alpha_t(k)&=b_k(x_t)\sum_{j=1}^{K}\alpha_{t-1}(j)A_{jk},\\
 p(z_t=k\mid x_{1:t})&=\frac{\alpha_t(k)}{\sum_j\alpha_t(j)}.
\end{align}
最后一时刻的前向量求和即为$p(x_{1:T})$。每步对$K^2$个状态对计算，总复杂度为$O(TK^2)$。长序列中概率连乘会下溢，应逐步缩放或在对数域使用 log-sum-exp，并保留缩放因子以恢复序列似然。

后向量$\beta_t(j)=p(x_{t+1:T}\mid z_t=j)$满足
\begin{equation}
 \beta_T(j)=1,\qquad
 \beta_t(j)=\sum_{k=1}^{K}A_{jk}b_k(x_{t+1})\beta_{t+1}(k).
\end{equation}
于是平滑概率为$\gamma_t(k)=\alpha_t(k)\beta_t(k)/p(x_{1:T})$。前向量汇总过去证据，后向量汇总未来证据，这也给出了链式消息传递的具体形式。

Viterbi 算法将前向递推中的求和替换为最大值。在对数域中，
\begin{equation}
 \delta_t(k)=\log b_k(x_t)+\max_j\{\delta_{t-1}(j)+\log A_{jk}\},
 \qquad \delta_1(k)=\log\pi_k+\log b_k(x_1).
\end{equation}
保存每步最大值对应的前驱，便可从终点回溯路径。逐时刻选择边缘概率最大的状态与 Viterbi 路径通常不同，前者优化逐点分类，后者优化整条路径的联合概率。在存在禁止转移时，逐点选择甚至可能产生不合法路径。

\paragraph{前向后向递推的推导与稳定实现。}
固定当前状态$k$，对前驱$j$边缘化：
\begin{align}
 p(x_{1:t},z_t=k)
 &=\sum_jp(x_t\mid z_t=k)p(z_t=k\mid z_{t-1}=j)
             p(x_{1:t-1},z_{t-1}=j)\\
 &=b_k(x_t)\sum_jA_{jk}\alpha_{t-1}(j).
\end{align}
条件独立允许把观测因子提到求和外。相似地，从$z_{t+1}$边缘化得到后向递推；给定$z_t$后过去和未来的观测分离，因此$\alpha_t(k)\beta_t(k)$就是$p(x_{1:T},z_t=k)$。求和表示保留所有可能路径，Viterbi的最大值则丢弃每个终点下非最优前缀；链结构保证最优前缀可扩展为最优整段路径。

数值上令$h_{t-1}(j)=p(z_{t-1}=j\mid x_{1:t-1})$，计算
\begin{equation}
 u_t(k)=b_k(x_t)\sum_jh_{t-1}(j)A_{jk},\quad
 c_t=\sum_ku_t(k),\quad h_t(k)=u_t(k)/c_t.
\end{equation}
第一步用$u_1(k)=\pi_kb_k(x_1)$。于是$c_t=p(x_t\mid x_{1:t-1})$且$\log p(x_{1:T})=\sum_t\log c_t$；缩放后的后向递推对应除以$c_{t+1}$。若$c_t=0$，是模型认为证据不可能而非普通下溢，应检查零转移与观测支持。对数域的$\log\sum_j e^{a_j}=m+\log\sum_je^{a_j-m}$（$m=\max_ja_j$）提供另一种稳定计算。

\paragraph{从隐状态推断到参数学习。}
隐状态未标注时，Baum--Welch 算法以 Expectation Maximization (EM) 交替估计状态的后验占用和模型参数。E 步计算$\gamma_t(j)$以及相邻状态的后验概率
\begin{equation}
 \xi_t(j,k)=\frac{\alpha_t(j)A_{jk}b_k(x_{t+1})\beta_{t+1}(k)}{p(x_{1:T})}.
\end{equation}
对一条序列，M 步的转移更新为
\begin{equation}
 A_{jk}^{\mathrm{new}}=
 \frac{\sum_{t=1}^{T-1}\xi_t(j,k)}{\sum_{t=1}^{T-1}\gamma_t(j)}.
\end{equation}
初始分布更新为$\pi_k^{\mathrm{new}}=\gamma_1(k)$，观测分布则按$\gamma_t(k)$进行加权估计。多条独立运行序列需要分别推断后再汇总计数，不能把不同设备的尾部与头部连接成一次状态转移。EM 依赖初始化，只保证在精确 E、M 步下不降低训练似然，不能保证找到全局最优参数。

\paragraph{EM为什么得到期望计数比。}
取$q(z)=p(z\mid x,\theta^{\rm old})$，由Jensen不等式，
\begin{equation}
 \log p_\theta(x)\geq\E_q\log p_\theta(x,z)-\E_q\log q(z).
\end{equation}
在旧参数处等号成立；M步保持$q$固定并提高右端，便不降低观测似然。转移相关部分是$\sum_{jk}N_{jk}\log A_{jk}$，其中$N_{jk}=\sum_{t<T}\xi_t(j,k)$。对每行加入约束乘子$\lambda_j(\sum_kA_{jk}-1)$，一阶条件$N_{jk}/A_{jk}+\lambda_j=0$与行和为一给出$A_{jk}=N_{jk}/\sum_kN_{jk}$。又因$\sum_k\xi_t(j,k)=\gamma_t(j)$，便得到前述更新。

若$x_t\in\R^d$且$b_k=\mathcal N(\mu_k,\Sigma_k)$，高斯对数密度的加权一阶条件给出
\begin{equation}
 \mu_k^{\rm new}=\frac{\sum_t\gamma_t(k)x_t}{N_k},\qquad
 \Sigma_k^{\rm new}=\frac{\sum_t\gamma_t(k)(x_t-\mu_k^{\rm new})(x_t-\mu_k^{\rm new})^{\mathsf T}}{N_k},
 \quad N_k=\sum_t\gamma_t(k).
\end{equation}
这里$\mu_k\in\R^d$、$\Sigma_k\in\R^{d\times d}$。若状态从未被占用，分母为零，参数不可由本轮数据识别；协方差也可能退化，需预先规定先验、协方差结构或数值下界，而非把无定义更新当作成功训练。约束或近似M步需要实际检查目标是否改善。

若状态转移已经观测到，可以为$A_{j,:}$设置 Dirichlet 先验，其后验参数等于先验参数加转移计数。状态未知时，转移计数本身也不确定。EM 的期望计数可用于点估计，但将其代入 Dirichlet 分布并不自动得到精确的参数后验，完整贝叶斯学习需要联合处理状态和参数的不确定性。

标准齐次 HMM 的自转移还隐含持续时间假设。若$A_{kk}<1$，进入状态$k$后停留$d$步的概率为
\begin{equation}
 p(D=d\mid k)=A_{kk}^{d-1}(1-A_{kk}),\qquad d=1,2,\ldots.
\end{equation}
这是一种无记忆的几何分布。若退化持续时间明显集中在某个范围，增加状态数未必是合适修复，隐半马尔科夫模型可以显式描述持续时间。动态贝叶斯网络则允许每个时刻包含多个相互作用的隐变量。线性高斯状态空间模型把离散状态换成连续状态，在相应假设下可用卡尔曼滤波精确计算过滤分布，后续序列章将讨论其动态含义。

\section{从局部条件模型到条件随机场}
最大熵马尔科夫模型与条件随机场分别代表局部条件归一化和全局路径归一化\cite{mccallum2000memm,lafferty2001crf}。比较二者的关键是概率质量如何在状态转移间分配，而不是是否使用神经网络提取特征。
HMM 同时学习观测如何产生和状态如何演化。当目标仅是识别状态，而输入包含频谱、载荷变化和多个重叠窗口特征时，准确建立$p(x_t\mid z_t)$可能比分类本身更困难。最大熵马尔科夫模型（Maximum Entropy Markov Model, MEMM）直接对状态转移建模，以对数线性分类器给出
\begin{equation}
 p_\theta(z_t=k\mid z_{t-1}=j,x)=
 \frac{\exp\{\theta^{\mathsf T}f(j,k,x,t)\}}
 {\sum_{k'}\exp\{\theta^{\mathsf T}f(j,k',x,t)\}}.
\end{equation}
其中$f$为特征向量，$x$表示当前任务允许使用的观测范围。最大熵的名称来自特征期望约束下的条件熵最大化，其解具有上述指数族形式。令$z_0$为固定起始状态，序列条件概率为各步条件概率之积。对于完整标注序列，训练可转化为以真实前驱状态为条件的局部分类；预测时前驱状态未知，仍需进行链上推断。

\begin{figure}[htbp]
\centering
\begin{tikzpicture}[>=Latex,
 state/.style={circle,draw,minimum size=10mm},
 every node/.style={font=\small}]
 \node[draw,rounded corners,fill=blue!12,minimum width=9cm,minimum height=9mm] (x) at (5,1.8) {已给定的可用观测 $x$（不为其建立生成分布）};
 \node[state,fill=gray!15] (z0) at (0,0) {$z_0$};
 \foreach \t in {1,2,3,4}{\node[state] (z\t) at (2.4*\t,0) {$z_{\t}$};}
 \foreach \a/\b in {0/1,1/2,2/3,3/4}{\draw[->,thick] (z\a) -- (z\b);}
 \foreach \t in {1,2,3,4}{\draw[->,blue!65!black] (x.south -| z\t.north) -- (z\t.north);}
 \node[align=center] at (4.8,-1.15) {每一步固定前驱 $z_{t-1}$，只对后继 $z_t$ 归一化\\$\displaystyle \sum_k p_\theta(z_t=k\mid z_{t-1},x)=1$};
\end{tikzpicture}
\caption{MEMM 的条件有向结构。灰色节点为固定起始状态，观测同时参与各步局部分类；与 HMM 相比，模型不再使用由状态指向观测的生成边。}
\label{fig:bayes-memm-structure}
\end{figure}
图\ref{fig:bayes-memm-structure}中，每个待预测标签的输入来自前驱标签与已给定观测，因此整条路径概率是这些局部条件概率的乘积。顶部的$x$可以包含任务允许的多个时刻和重叠特征，并不意味着在线任务能够读取未来。横向箭头保留了标签顺序依赖，却没有取消每一步各自归一化的限制；这正是后文标签偏置的结构来源。

\paragraph{最大熵为何导出softmax。}
把前驱与可用输入合称上下文$c$，令经验上下文权重为$\hat p(c)$，特征$\phi(c,k)\in\R^p$。在归一化约束$\sum_kq(k\mid c)=1$及特征期望约束$\sum_{c,k}\hat p(c)q(k\mid c)\phi(c,k)=\bar\phi$下，最大化条件熵$-\sum_{c,k}\hat p(c)q(k\mid c)\log q(k\mid c)$。对$q(k\mid c)$求拉格朗日一阶条件得到
\begin{equation}
 -\log q(k\mid c)-1+\theta^{\mathsf T}\phi(c,k)+\lambda(c)=0.
\end{equation}
指数化后，归一化消去$\lambda(c)$，即得$q(k\mid c)=\exp(\theta^{\mathsf T}\phi(c,k))/\sum_{k'}\exp(\theta^{\mathsf T}\phi(c,k'))$。这要求约束可行；特征完全分离时有限最优权重未必存在，通常以正则化松弛精确矩约束。名称中的熵指这个构造原则，不表示模型输出必须高熵。

MEMM 可以使用相关、重叠的输入特征，不需要描述它们的联合分布，但每个前驱状态都单独归一化，这带来了标签偏置。若某个状态只有一个允许的后继，该转移的条件概率恒为一，无论对应特征得分多低，都不能借助该步降低路径概率。更一般地，后继较少或转移分布较集中的状态可能获得结构性优势。问题来自局部归一化对路径评分的约束，并非简单的类别数量不平衡。

线性链条件随机场（CRF）保留相邻标签依赖，却把归一化放到整条路径上\cite{murphy2022pml}。定义路径得分
\begin{align}
 F_\theta(z,x)&=\sum_{t=1}^{T}\theta^{\mathsf T}f(z_{t-1},z_t,x,t),\\
 p_\theta(z\mid x)&=\frac{\exp F_\theta(z,x)}{Z_\theta(x)},\\
 Z_\theta(x)&=\sum_{z'_{1:T}}\exp F_\theta(z',x).
\end{align}
\begin{figure}[htbp]
\centering
\begin{tikzpicture}[
 state/.style={circle,draw,minimum size=9mm},
 factor/.style={rectangle,draw,fill=orange!25,minimum size=7mm,inner sep=1pt},
 every node/.style={font=\small}]
 \node[draw,rounded corners,fill=blue!12,minimum width=10cm,minimum height=9mm] (x) at (5,1.8) {给定观测 $x$：各局部势函数共享的条件};
 \node[state,fill=gray!15] (z0) at (0,0) {$z_0$};
 \foreach \t in {1,2,3,4}{
  \node[state] (z\t) at (2.5*\t,0) {$z_{\t}$};
  \node[factor] (f\t) at (2.5*\t-1.25,0) {$\psi_{\t}$};
  \draw[thick] (f\t) -- (z\t);
  \draw[dashed,blue!65!black] (x.south -| f\t.north) -- (f\t.north);
 }
 \foreach \a/\b in {0/1,1/2,2/3,3/4}{\draw[thick] (z\a) -- (f\b);}
 \node[align=center] at (5,-1.2) {$\psi_t=\exp\{\theta^{\mathsf T}f(z_{t-1},z_t,x,t)\}$\\局部势相乘，再用同一个 $Z_\theta(x)$ 对整条路径归一化};
\end{tikzpicture}
\caption{线性链 CRF 的因子图。圆形表示标签，方形表示相邻标签的势函数；实线无方向，表示变量参与因子。蓝色虚线仅表示势函数依赖已给定观测，不是标签图中的随机变量边或生成箭头。}
\label{fig:bayes-crf-structure}
\end{figure}
图\ref{fig:bayes-crf-structure}把条件分布写成可逐项检查的乘积：每个方形读取相邻两个标签及观测，输出一个非负相容度。固定$z_0$后，第一个因子实际只涉及一个待预测标签。去掉因子节点、将共享因子的标签直接相连，就得到标签的无向链。与 MEMM 不同，$\psi_t$不要求对后继标签求和为一，所有因子共同决定路径得分，最后才由一个全局配分函数归一化。

配分函数$Z_\theta(x)$比较所有允许路径的得分。一个只有唯一后继的状态仍然可以获得较低的路径得分，因此消除了 MEMM 局部归一化造成的上述限制。全局归一化不意味着自动解决标注噪声、分布偏移或特征不足。

对一条标注序列$(x,z)$，带二次正则的负对数条件似然为
\begin{equation}
 \mathcal L(\theta)=-F_\theta(z,x)+\log Z_\theta(x)
 +\frac{\lambda}{2}\|\theta\|_2^2.
\end{equation}
其梯度等于模型期望的特征计数减去真实路径的特征计数，再加$\lambda\theta$。这个形式解释了训练为何需要边缘概率：只找出当前最优路径不足以计算配分函数的梯度。在线性链、有限状态和相邻标签因子下，前向后向算法可在$O(TK^2)$时间内计算配分函数和边缘概率，Viterbi 则计算最优路径。

\paragraph{配分函数、边缘概率与梯度的完整联系。}
令$\psi_t(j,k)=\exp\{\theta^{\mathsf T}f(j,k,x,t)\}$。固定起点$z_0$，设置$a_0(z_0)=1$、其余为零，递推$a_t(k)=\sum_ja_{t-1}(j)\psi_t(j,k)$，则$Z=\sum_ka_T(k)$。后向量$b_T(k)=1$、$b_{t-1}(j)=\sum_k\psi_t(j,k)b_t(k)$给出边缘
\begin{equation}
 p_\theta(z_{t-1}=j,z_t=k\mid x)
 =\frac{a_{t-1}(j)\psi_t(j,k)b_t(k)}{Z}.
\end{equation}
这些$a,b$是未归一化消息，不是HMM的观测似然参数。定义总特征
\begin{equation}
 \Phi(z,x)=\sum_tf(z_{t-1},z_t,x,t)\in\R^p.
\end{equation}
直接对有限和求导得
\begin{align}
 \nabla_\theta\log Z
 &=\frac{\sum_z e^{\theta^{\mathsf T}\Phi(z,x)}\Phi(z,x)}{Z}
 =\E_{p_\theta(z\mid x)}\Phi(z,x),\\
 \nabla_\theta\mathcal L
 &=-\Phi(z^{\rm obs},x)+\sum_{t,j,k}p_\theta(j,k\mid x)f(j,k,x,t)+\lambda\theta.
\end{align}
再求导得到$\nabla^2\mathcal L=\operatorname{Cov}_{p_\theta}(\Phi)+\lambda I_p$。因此固定特征的线性CRF目标凸，$\lambda>0$时严格凸；若特征由待训练神经网络产生，这一参数凸性便不再成立。禁止路径可赋零势，但至少必须存在一条合法路径，否则$Z=0$使条件分布无定义。

神经网络可以提供 CRF 的局部观测得分。例如，令$s_t(k;x)$为编码器输出，$B_{jk}$为可学习转移得分，则
\begin{equation}
 F_\theta(z,x)=\sum_{t=1}^{T}s_t(z_t;x)+\sum_{t=2}^{T}B_{z_{t-1},z_t}.
\end{equation}
与 HMM 的$A_{jk}$不同，$B_{jk}$是未归一化得分，不要求非负或行和为一。Bidirectional Long Short-Term Memory (BiLSTM)--CRF 和 Transformer--CRF 将特征提取交给神经网络，将标签关系交给结构化输出层。双向编码器适合完整记录的离线标注；在线预警必须限制输入和推断范围，或者明确等待未来观测所引入的延迟。即使局部特征只读取过去，对整段序列进行平滑仍可能利用后续证据。

\section{马尔科夫随机场与空间一致性}
图像中的局部相互作用可用 Gibbs 分布和马尔科夫随机场描述，Geman 与 Geman 将这一建模方式用于图像恢复\cite{geman1984mrf}。本节据此解释邻域势函数与能量最小化，但不把平滑先验当作真实边界的替代。
时间链只是一种邻接关系。在缺陷图像或设备网络中，每个变量可能有多个邻居，马尔科夫随机场（Markov Random Field, MRF）用无向图描述这种依赖。对严格正的分布，在相应马尔科夫性质下可以按图中的团分解为
\begin{equation}
 p(z)=\frac{1}{Z}\prod_{C\in\mathcal C}\psi_C(z_C)
 =\frac{1}{Z}\exp\{-E(z)\}.
\end{equation}
\begin{figure}[htbp]
\centering
\begin{tikzpicture}[
 site/.style={circle,draw,minimum size=11mm,fill=white},
 every node/.style={font=\small}]
 \foreach \r in {1,2,3}{
  \foreach \c in {1,2,3}{\node[site] (v\r\c) at (1.8*\c,-1.5*\r) {$z_{\r\c}$};}
 }
 \foreach \r in {1,2,3}{\foreach \a/\b in {1/2,2/3}{\draw[thick] (v\r\a) -- (v\r\b);}}
 \foreach \c in {1,2,3}{\foreach \a/\b in {1/2,2/3}{\draw[thick] (v\a\c) -- (v\b\c);}}
 \node[site,fill=orange!30] at (v22) {$z_{22}$};
 \foreach \n in {12,21,23,32}{\node[site,fill=green!15] at (v\n) {$z_{\n}$};}
 \node[align=left,text width=5.2cm,anchor=west] at (6.3,-3) {橙色：当前标签 $z_{22}$\\[3pt]绿色：四邻域 $N(22)$\\[3pt]白色：其余标签\\[6pt]给定绿色邻居后，中心标签与其余标签条件独立。};
\end{tikzpicture}
\caption{四邻域网格上的成对 MRF。每个节点对应一个像素或区域标签，每条无向边对应一个二元势；不存在箭头，并不表示各标签独立，而是表示相互作用不采用有向生成分解。}
\label{fig:bayes-mrf-structure}
\end{figure}
图\ref{fig:bayes-mrf-structure}展示的是一种具体的网格结构，并非所有 MRF 都必须是网格。此处$z_{rc}$的两个下标分别表示行和列。局部马尔科夫性质可写为
\begin{equation}
 p(z_{22}\mid z_{\setminus 22})
 =p(z_{22}\mid z_{12},z_{21},z_{23},z_{32}),
\end{equation}
其中$z_{\setminus 22}$表示除中心外的全部标签。对于图示的一元和二元势分解，固定其他标签后，与$z_{22}$无关的因子在条件概率归一化时相消，因而
\begin{equation}
 p(z_{22}=k\mid z_{\setminus 22})
 =\frac{\psi_{22}(k)\prod_{j\in N(22)}\psi_{22,j}(k,z_j)}
 {\sum_{k'}\psi_{22}(k')\prod_{j\in N(22)}\psi_{22,j}(k',z_j)}.
\end{equation}
这也解释了 Gibbs 采样为何只需读取当前节点的邻居。不过，局部更新只访问邻域，并不意味着联合分布能够分解成相互独立的节点分布。若这些势还依赖给定图像$x$，图中的标签邻接关系保持不变，得到的就是后文讨论的网格 CRF。

团$C$内的节点两两相连，势函数$\psi_C$表达局部相容程度，并不是归一化概率。能量$E(z)=-\sum_C\log\psi_C(z_C)$越低，配置概率越高。配分函数$Z$对全部配置求和，保证总概率为一。若使用零势函数表达硬约束，则不能直接套用依赖严格正性的等价结论。

在给定图像$x$的二元缺陷标注中，可以设置
\begin{equation}
 E(z;x)=\sum_i U_i(z_i;x)+\lambda\sum_{(i,j)\in\mathcal E}
 w_{ij}(x)\mathbb 1\{z_i\ne z_j\},\qquad \lambda,w_{ij}(x)\geq0.
\end{equation}
一元项$U_i$衡量局部像素证据，二元 Potts 项惩罚邻居标签不一致。相邻区域越相似，$w_{ij}$可以越大；跨越明显边界时降低权重，可减轻过度平滑。若直接将分类概率转换成一元负对数得分，这构成条件模型的一种参数化，并不自动成为生成似然。

此时$p(z\mid x)\propto\exp\{-E(z;x)\}$也是一个 CRF。MRF 强调无向马尔科夫依赖，CRF 强调给定观测后的结构化条件分布，两者不是互斥的模型名单。线性链 CRF 是其中计算特别方便的一类，二维网格 CRF 通常不能直接沿用链式动态规划。

能量中的平滑项可以移除孤立误报，也可能抹掉真实细裂纹，因此应依据缺陷宽度和边界误差选择强度。一般含环图的精确推断可能十分昂贵，变量消元的复杂度随树宽增长。置信传播在树上精确，在含环图上通常只是近似且不保证收敛。均值场以较简单的分布逼近联合依赖，Gibbs 采样则逐变量更新。对于二元、次模的成对能量，图割可求全局最小能量，但这一结论不能推广到任意多类别势函数或边缘概率计算。

\begin{table}[htbp]
\centering
\caption{链式与无向图模型的建模对象和计算边界。}
\label{tab:bayes-graph-models}
\small
\begin{tabular}{p{1.3cm}p{3.1cm}p{3.3cm}p{4.2cm}}
\hline
模型 & 建模对象 & 归一化方式 & 关键限制 \\
\hline
HMM & 联合分布$p(x,z)$ & 转移与观测局部归一化 & 观测假设和状态持续时间需检验 \\
MEMM & 条件分布$p(z\mid x)$ & 每个前驱状态局部归一化 & 标签偏置影响路径评分 \\
链式 CRF & 条件分布$p(z\mid x)$ & 所有标签路径全局归一化 & 训练需计算模型期望 \\
MRF & 无向图上的联合分布 & 全部配置全局归一化 & 一般图中配分函数和推断困难 \\
\hline
\end{tabular}
\end{table}

表\ref{tab:bayes-graph-models}按建模对象、归一化方式和推断边界比较图模型；应先决定需要联合生成还是条件预测，再判断链式动态规划是否适用。
\paragraph{MRF的局部更新与均值场。}
对固定其他标签$z_{-i}$，不包含$i$的能量项在条件概率的分子分母中相消，因此
\begin{equation}
 p(z_i=k\mid z_{-i},x)\propto
 \exp\left\{-U_i(k;x)-\lambda\sum_{j:(i,j)\in\mathcal E}w_{ij}(x)\mathbb 1\{k\ne z_j\}\right\}.
\end{equation}
逐点从此分布抽样就是Gibbs更新；若改为取最小局部能量，则变为坐标下降，只能保证能量不增而非全局最优。有限状态、适当遍历性和非周期性下Gibbs链以目标分布为平稳分布，强耦合时仍可能混合极慢。

若用乘积分布$q(z)=\prod_iq_i(z_i)$近似，最小化$\KL(q\|p)$等价于最小化以下目标：
\begin{equation}
 \E_qE+\sum_i\sum_kq_i(k)\log q_i(k).
\end{equation}
固定其他因子并对$q_i$加归一化乘子求导，可得
\begin{equation}
 q_i(k)\propto\exp\{-\E_{q_{-i}}E(k,z_{-i};x)\}
 \propto\exp\{-U_i(k;x)-\lambda\sum_jw_{ij}(1-q_j(k))\}.
\end{equation}
这以邻居概率代替确定标签，但忽略联合相关性；逐坐标精确更新改善变分目标，却不保证找到全局最佳近似。二元成对能量可用图割的关键条件是$V_{ij}(0,0)+V_{ij}(1,1)\leq V_{ij}(0,1)+V_{ij}(1,0)$，非负Potts项满足它。图割解决最大概率配置，不提供配分函数或后验区间。

\paragraph{把图结构带回维护任务。}
假设要从一次运行记录中识别正常、退化与故障阶段，可以先用独立分类器建立逐点基线。若预测频繁跳变，HMM 用转移持续性与观测模型共同解释状态；若观测特征复杂而状态标注充分，MEMM 和 CRF 可直接学习条件边界。若任务转为表面缺陷的空间连通区域，则链结构不再适用，需要结合图像邻接关系定义条件随机场。模型变化由输出依赖决定，而非由名称的新旧决定。

这些比较必须使用相同的可用观测、标签定义和设备划分。逐点准确率之外，还应检查状态切换延迟、阶段持续时间误差和关键故障段漏检。离线平滑结果只能与离线结果比较，在线过滤应另报实际延迟。若由隐状态概率推导未来故障风险，还需要定义故障事件与状态的关系，并验证预测时间窗，不能把“当前处于退化状态”的概率直接解释成“未来二十四小时故障”的概率。

上述动态规划在参数给定时可以精确计算链式隐状态分布，但这不意味着参数已经确定。小样本设备的转移矩阵、观测噪声和条件特征权重仍有估计误差。接下来回到参数层，讨论哪些模型可以解析更新，哪些模型需要近似后验。

\section{从共轭模型到近似推断}
\paragraph{共轭模型与近似推断。}
若先验与似然属于共轭族，后验可解析求得。例如，伯努利观测的成功概率$p$取Beta先验$p\sim\mathrm{Beta}(a,b)$，观察到$k$次成功、$n-k$次失败后有
\begin{equation}
 p\mid\mathcal{D}\sim\mathrm{Beta}(a+k,b+n-k).
\end{equation}
后验均值为
\begin{equation}
 \mathbb{E}[p\mid\mathcal D]=\frac{a+k}{a+b+n}.\label{eq:bayes-beta-mean}
\end{equation}
式\eqref{eq:bayes-beta-mean}将先验均值与样本比例按先验强度和观测数量加权。$a+b$越大，先验影响越强。用于故障率时，这个更新假定各次观测在给定$p$后独立且共享同一事件概率，不能直接把大量重叠时间窗视为独立试验。

\paragraph{Beta更新的归一化与预测。}
对$a,b>0$，先验密度为$p^{a-1}(1-p)^{b-1}/B(a,b)$，其中$B(a,b)=\int_0^1u^{a-1}(1-u)^{b-1}\,du$。给定独立伯努利序列，似然正比于$p^k(1-p)^{n-k}$，相乘后指数为$a+k-1$和$b+n-k-1$，用$B(a+k,b+n-k)$归一化便得到后验。若观测对象只是计数，似然还有$\binom nk$，它不依赖$p$，不改变后验。

记$A=a+k$、$B=b+n-k$（此处大写$B$为参数，不是Beta函数），由归一化积分之比得到
\begin{equation}
 \E[p\mid\mathcal D]=\frac{A}{A+B},\quad
 \Var(p\mid\mathcal D)=\frac{AB}{(A+B)^2(A+B+1)},\quad
 P(Y_*=1\mid\mathcal D)=\E[p\mid\mathcal D].
\end{equation}
其中$\E p=\frac{a+b}{a+b+n}\frac a{a+b}+\frac n{a+b+n}\frac kn$（$n>0$），显式给出收缩权重。未来$m$次试验的成功数$K_*$满足
\begin{equation}
 P(K_*=r\mid\mathcal D)=\binom mr
 \frac{B(A+r,B+m-r)}{B(A,B)},\quad 0\leq r\leq m.
\end{equation}
这是先对二项分布的未知$p$积分所得，不等同于把后验均值代入二项分布：积分后未来试验通过共享$p$相关，从而保留参数不确定性。

共轭关系也适用于连续参数。考虑$y=X\beta+\varepsilon$，其中$\varepsilon\sim\mathcal N(0,\sigma^2I)$，先验$\beta\sim\mathcal N(m_0,S_0)$。当$\sigma^2$已知时，贝叶斯线性回归的后验仍为高斯分布，其精度矩阵和均值为
\begin{align}
 S_n^{-1}&=S_0^{-1}+\sigma^{-2}X^{\mathsf T}X,\\
 m_n&=S_n\left(S_0^{-1}m_0+\sigma^{-2}X^{\mathsf T}y\right).
\end{align}
先验精度与数据精度相加，说明收缩先验如何稳定共线特征下的小样本回归。对新输入$x_*$，预测均值为$x_*^{\mathsf T}m_n$，预测方差为$\sigma^2+x_*^{\mathsf T}S_nx_*$，分别包含观测噪声和参数不确定性。

\paragraph{高斯后验的配方步骤。}
令$X\in\R^{n\times d}$、$\beta,m_0\in\R^d$、$S_0\in\R^{d\times d}$正定，$y\in\R^n$，且$\sigma^2>0$已知。负对数联合密度中依赖$\beta$的部分为
\begin{align}
 &\frac12(\beta-m_0)^{\mathsf T}S_0^{-1}(\beta-m_0)
   +\frac1{2\sigma^2}(y-X\beta)^{\mathsf T}(y-X\beta)\\
 &=\frac12\beta^{\mathsf T}\Lambda\beta-\beta^{\mathsf T}h+\text{常数}\\
 &=\frac12(\beta-\Lambda^{-1}h)^{\mathsf T}\Lambda
   (\beta-\Lambda^{-1}h)+\text{常数},
\end{align}
其中$\Lambda=S_0^{-1}+\sigma^{-2}X^{\mathsf T}X$、$h=S_0^{-1}m_0+\sigma^{-2}X^{\mathsf T}y$。识别高斯指数可得$S_n=\Lambda^{-1}$、$m_n=S_nh$。正定先验使$X$秩不足时后验仍适定；实际通过线性方程求解，不必显式求逆。

新观测为$Y_*=x_*^{\mathsf T}\beta+\varepsilon_*$，其中$x_*\in\R^d$且新噪声独立，因此线性高斯和仍为高斯，均值$x_*^{\mathsf T}m_n$、方差$x_*^{\mathsf T}S_nx_*+\sigma^2$。潜在均值区间只含第一项，观测预测区间包含两项。若噪声尺度未知并为其指定共轭先验，积分后常得到Student分布，而非直接沿用已知$\sigma^2$的高斯结论。

复杂深度模型通常没有解析后验，可采用最大后验估计、Laplace近似、变分推断或Markov Chain Monte Carlo (MCMC)。变分推断选择易处理的$q_{\phi}(\bm{\theta})$逼近后验，并最小化
\begin{equation}
 \KL\left(q_{\phi}(\bm{\theta})\,\|\,p(\bm{\theta}\mid\mathcal{D})\right).
\end{equation}
通过证据下界（Evidence Lower Bound，ELBO）可写成
\begin{equation}
 \log p(\mathcal{D})\geq \mathbb{E}_{q_{\phi}}[\log p(\mathcal{D}\mid\bm{\theta})]-\KL(q_{\phi}(\bm{\theta})\|p(\bm{\theta})).\label{eq:bayes-elbo}
\end{equation}
式\eqref{eq:bayes-elbo}右端的第一项鼓励解释数据，第二项控制近似后验偏离先验的程度。

这一优化目标来自恒等式
\begin{equation}
 \KL(q_\phi(\theta)\|p(\theta\mid\mathcal D))
 =\log p(\mathcal D)-\operatorname{ELBO}(\phi).
\end{equation}
证据对$\phi$为常数，因此最大化 ELBO 等价于最小化所选方向的 KL 散度。均值场近似$q(\theta)=\prod_jq_j(\theta_j)$计算方便，却忽略参数相关性，可能低估后验方差。结构化变分分布保留部分依赖，代价是更复杂的计算。

若变分变量$z\sim\mathcal N(\mu_\phi,\sigma_\phi^2)$，可写成$z=\mu_\phi+\sigma_\phi\epsilon$，其中$\epsilon\sim\mathcal N(0,1)$与$\phi$无关。重参数化把随机性移到外部噪声，使 ELBO 的梯度能够通过可微变换估计。离散隐状态则未必具有这种可微表示，应优先检查是否能像 HMM 那样精确求和，再考虑松弛或得分函数估计。

\paragraph{ELBO恒等式与可计算梯度。}
设$q$在联合密度有支持的区域内，且所需积分有限。把Bayes公式代入KL定义，逐项展开为
\begin{align}
 \KL(q\|p(\theta\mid\mathcal D))
 &=\E_q\log q-\E_q\log p(\mathcal D,\theta)+\log p(\mathcal D),\\
 \operatorname{ELBO}(q)
 &=\E_q\log p(\mathcal D\mid\theta)+\E_q\log p(\theta)-\E_q\log q.
\end{align}
KL非负给出下界，等号只在$q$等于真实后验（几乎处处）时成立。证据最大并不意味着近似最准确：若不同模型的变分族表达能力不同，下界差距还混有近似误差。

对均值场$q(\theta)=\prod_jq_j(\theta_j)$，固定其他因子，把ELBO视为$q_j$的函数并加归一化乘子，可得
\begin{equation}
 \log q_j^*(\theta_j)=\E_{q_{-j}}\log p(\mathcal D,\theta)+\text{归一化常数}.
\end{equation}
这给出坐标上升更新的来源；期望或归一化未必能解析算出。若$\theta=\mu+s\odot\epsilon$、$\epsilon\sim\mathcal N(0,I_d)$、$s_j>0$，在可交换微分与积分的可积条件下，
\begin{equation}
 \nabla_\mu\E_\epsilon h(\mu+s\odot\epsilon)
 =\E_\epsilon\nabla_\theta h(\theta),\qquad
 \partial_{s_j}\E_\epsilon h(\theta)=\E_\epsilon[\partial_{\theta_j}h(\theta)\epsilon_j].
\end{equation}
正态$q$的熵为$\frac12\sum_j\log(2\pi e s_j^2)$，其对$s_j$的导数为$1/s_j$，与联合对数密度梯度合成完整ELBO梯度。可用$s_j=\exp(\rho_j)$保证正值。采样梯度只近似当前变分目标的梯度，不能消除变分族的结构偏差。

\paragraph{贝叶斯更新的几何直观。}
贝叶斯公式不仅是概率比例式，也可以理解为在函数空间中重新分配权重。先验把概率质量放在一组可能的参数上，似然根据参数对观测数据的解释能力重新加权，后验再进行归一化。数据量较小时，先验仍然有明显影响；数据量增加后，似然通常占据主导，但这依赖模型假设没有严重错配。

\begin{figure}[htbp]
\centering
\begin{tikzpicture}[>=Latex, node distance=1.8cm, every node/.style={font=\small}]
 \node[draw, rounded corners, fill=blue!8, minimum width=2.6cm, minimum height=1cm] (prior) {先验 $p(\theta)$};
 \node[draw, rounded corners, fill=orange!12, right=of prior, minimum width=2.8cm, minimum height=1cm] (like) {似然 $p(\mathcal D\mid\theta)$};
 \node[draw, rounded corners, fill=green!12, right=of like, minimum width=2.8cm, minimum height=1cm] (post) {后验 $p(\theta\mid\mathcal D)$};
 \draw[->, thick] (prior) -- node[above]{结合} (like);
 \draw[->, thick] (like) -- node[above]{归一化} (post);
 \draw[->, dashed] ($(prior.south)!0.5!(like.south)$) .. controls +(0,-1.0) and +(0,-1.0) .. node[below]{模型假设} ($(like.south)!0.5!(post.south)$);
\end{tikzpicture}
\caption{贝叶斯更新把先验结构与数据证据结合为后验。}
\label{fig:bayes-update-components}
\end{figure}
图\ref{fig:bayes-update-components}将先验与数据证据汇入后验，强调后验结论同时依赖结构假设和观测；检验预测可信度时不能只检查数据拟合。

\section{模型比较与不确定性计算}
\paragraph{从 Maximum A Posteriori (MAP) 到正则化。}
最大后验估计最大化
\begin{equation}
 \hat\theta_{MAP}=\argmax_\theta \log p(\mathcal D\mid\theta)+\log p(\theta).
\end{equation}
以高斯先验$p(\theta)\propto\exp(-\lambda\|\theta\|_2^2/2)$为例，负对数后验等价于经验损失加上$\lambda\|\theta\|_2^2/2$。因此正则化并不是一个与概率建模无关的工程技巧，而是在特定先验下的推断结果。需要注意的是，正则化系数、噪声方差和先验方差之间存在耦合，不能脱离似然尺度解释$\lambda$。

\paragraph{边际似然与模型比较。}
边际似然
\begin{equation}
 p(\mathcal D)=\int p(\mathcal D\mid\theta)p(\theta)\,d\theta
\end{equation}
同时惩罚拟合不足和模型过度灵活。它与只比较训练集似然不同，因为参数空间中大量不能解释数据的区域会降低平均证据。实际计算边际似然通常困难，可使用 ELBO、Laplace 近似或信息准则近似。模型比较时还应明确先验是否改变，因为不同先验下的边际似然不具有无条件可比性。

\paragraph{Laplace 近似。}
在后验众数$\hat\theta$附近，对负对数后验$U(\theta)=-\log p(\theta\mid\mathcal D)$做二阶 Taylor 展开：
\begin{equation}
 U(\theta)\approx U(\hat\theta)+\frac12(\theta-\hat\theta)^{\mathsf T}H(\theta-\hat\theta),
 \qquad H=\nabla^2U(\hat\theta).
\end{equation}
于是后验近似为$\mathcal N(\hat\theta,H^{-1})$。该近似在后验近似单峰且局部二次时有效；深度网络存在对称性、多峰和平坦方向时，单个高斯可能严重低估不确定性。

\paragraph{从局部二次式到归一化常数。}
若$\hat\theta\in\R^d$为内部驻点且$H$正定，线性项因梯度为零而消失。指数化二阶式，并用高斯积分
\begin{equation}
 \int\exp\{-\tfrac12u^{\mathsf T}Hu\}\,du=(2\pi)^{d/2}|H|^{-1/2}
\end{equation}
归一化，便得到协方差$H^{-1}$。若从未归一化密度$\tilde p(\theta)=p(\mathcal D\mid\theta)p(\theta)$出发，同样推得
\begin{equation}
 \log p(\mathcal D)\approx\log\tilde p(\hat\theta)
 +\frac d2\log(2\pi)-\frac12\log|H|.
\end{equation}
它同时考虑峰高与峰附近体积，而不是只看最大似然。边界众数、奇异Hessian、多峰或不适当先验会使这一步失效；加入数值阻尼可改善线性代数，却不是证明原后验本来就是该高斯。

\paragraph{MCMC 的接受率直观。}
Metropolis--Hastings 从当前状态$\theta$提出候选$\theta'$，以
\begin{equation}
 \alpha=\min\left\{1,\frac{p(\theta'\mid\mathcal D)q(\theta\mid\theta')}{p(\theta\mid\mathcal D)q(\theta'\mid\theta)}\right\}
\end{equation}
的概率接受。候选分布步长过小会导致样本高度相关，步长过大则接受率过低。实际诊断不能只看链是否“运行完成”，还应检查 burn-in、有效样本量、自相关和多条链之间的混合。

\begin{table}[htbp]
\centering
\caption{贝叶斯推断方法的工程取舍}
\label{tab:bayes-inference-methods}
\scriptsize
\setlength{\tabcolsep}{3pt}
\resizebox{\linewidth}{!}{%
\begin{tabular}{p{2.5cm}p{3.6cm}p{3.6cm}p{4.0cm}}
\hline
方法 & 近似对象 & 优势 & 主要风险与适用场景 \\
\hline
MAP & 后验众数 & 成本低、易扩展、适合在线更新 & 忽略参数不确定性，适合作为基线 \\
Laplace & 众数附近的高斯曲率 & 能快速产生局部区间 & 多峰、对称和非二次后验会导致低估 \\
变分推断 & 可处理的$q_\phi(\theta)$ & 可用随机梯度扩展到大模型 & 均值场假设常低估相关性和方差 \\
MCMC & 后验样本链 & 理论上逼近目标后验 & 计算昂贵，必须检查混合和有效样本量 \\
\hline
\end{tabular}}
\end{table}

表\ref{tab:bayes-inference-methods}区分局部近似、变分优化与采样的计算取舍；选用方法时应把其近似误差和诊断要求与后验形状共同考虑。
\paragraph{接受率为何保持目标分布。}
记目标密度为$\pi(\theta)$。对不同状态$\theta,\theta'$，Metropolis--Hastings (MH)的双向概率流满足
\begin{align}
 \pi(\theta)q(\theta'\mid\theta)\alpha(\theta,\theta')
 &=\min\{\pi(\theta)q(\theta'\mid\theta),
          \pi(\theta')q(\theta\mid\theta')\}\\
 &=\pi(\theta')q(\theta\mid\theta')\alpha(\theta',\theta).
\end{align}
这就是细致平衡；拒绝时停留原位使转移核的总概率为一。对当前状态积分可验证$\pi$不变，且后验的未知证据常数在比值中消去。平稳性不等于有限步已收敛；还需要能到达目标支持、非周期性及相应遍历条件。对称提议可约去$q$，非对称提议遗漏该比值则会改变目标分布。

\section{校准、诊断与工业决策}
\paragraph{不确定性的分解与校准。}
预测方差可用全概率方差公式分解：
\begin{equation}
 \Var(y_*\mid x_*,\mathcal D)=\mathbb E_{p(\theta\mid\mathcal D)}[\Var(y_*\mid x_*,\theta)]
 +\Var_{p(\theta\mid\mathcal D)}[\mathbb E(y_*\mid x_*,\theta)].\label{eq:bayes-predictive-variance}
\end{equation}
式\eqref{eq:bayes-predictive-variance}右端的第一项主要对应观测噪声，第二项对应参数不确定性。模型给出的区间还必须经过覆盖率检查：预测百分之九十区间不应只在训练分布内覆盖百分之六十的真实值。对于工业决策，区间宽度与错误成本同样重要。

\paragraph{全方差公式的推导与估计。}
省略给定的$x_*,\mathcal D$，记$m_\theta=\E[Y_*\mid\theta]$、$m=\E_\theta m_\theta$，则$Y_*-m=(Y_*-m_\theta)+(m_\theta-m)$。平方后取期望，交叉项为零，因为$\E[Y_*-m_\theta\mid\theta]=0$，从而得到全方差公式。这个分解是在指定模型内的，不包含模型类别错设所遗漏的不确定性。

若有$S$个参数样本，并能计算条件均值$m_s$及方差$v_s$，令$\bar m=S^{-1}\sum_sm_s$，则经验混合分布的精确方差为
\begin{equation}
 \hat v=\frac1S\sum_sv_s+\frac1S\sum_s(m_s-\bar m)^2.
\end{equation}
若目的是无偏估计独立后验样本的均值间方差，第二项可改用$S-1$分母，但这时不再是该有限混合的精确方差。回归区间需从完整预测混合求分位数，不能仅对$m_s$取分位数后声称覆盖有噪声观测。对于重尾或多峰预测，$\bar m\pm1.96\sqrt{\hat v}$也不一定是95\%区间。

\paragraph{贝叶斯方法的常见误用。}
使用贝叶斯模型时，可以从结果倒查计算过程。先检查报告的概率与实际发生频率是否相符，再检查陌生工况下模型是否仍然过度自信。后验方差只反映当前模型表达的未知量，模型没有考虑到的故障机制可能根本不会使它增大。先验也应写明来源和尺度；如果近似算法尚未稳定，计算误差甚至可能大于想要估计的参数不确定性。

\paragraph{预测分布的蒙特卡洛近似。}
当后验积分无法解析计算时，可从后验得到$S$个样本$\theta^{(s)}$，使用
\begin{equation}
 p(y_*\mid x_*,\mathcal D)\approx\frac1S\sum_{s=1}^{S}p(y_*\mid x_*,\theta^{(s)}).
\end{equation}
注意先平均概率再取最大类别；平均 logits 或平均类别标签一般不等价。对回归，可由每个参数样本产生一条预测，再计算总方差和分位数区间。$S$太小会造成蒙特卡洛噪声，样本相关时有效样本量还会进一步下降。

\paragraph{MCMC 诊断不能省略。}
链的轨迹图可以检查是否进入稳定区域，自相关函数衡量样本冗余，有效样本量近似为
\begin{equation}
 ESS\approx\frac{S}{1+2\sum_{k\geq1}\rho_k}.
\end{equation}
多条不同初始状态的链还应比较$\hat R$等收敛指标。若链只在后验的一个模式附近游走，表面上的稳定并不代表全局混合良好。工业高维模型通常更偏向近似推断，但这不意味着可以跳过近似质量和预测覆盖率检查。

\paragraph{Effective Sample Size (ESS)是针对哪个估计量。}
对平稳链上的标量$h(\theta^{(s)})$，设其方差为$v_h$、滞后相关为$\rho_k$，展开样本均值的协方差和得
\begin{equation}
 \Var\left(\frac1S\sum_sh(\theta^{(s)})\right)
 =\frac{v_h}{S}\left[1+2\sum_{k=1}^{S-1}(1-k/S)\rho_k\right].
\end{equation}
当相关可求和且链足够长时，匹配独立样本均值的方差便得到前述ESS。不同参数、尾部概率和维护成本函数有不同的自相关，因此不存在适用于所有结论的单一ESS。短链直接相加所有经验相关会很不稳定，需截断或稳定的谱估计；尚未进入平稳区或未跨模式时，这个局部诊断不能证明准确性。

\paragraph{先验敏感性分析。}
先验选择应通过敏感性分析验证。对故障率、噪声尺度或参数收缩强度使用几组合理先验，比较后验预测和决策是否稳定。若结果对先验极其敏感，说明数据不足、模型可识别性弱，或先验承担了过多结论。此时应增加数据、改进测量或明确报告先验依赖，而不是只选择产生最好结果的先验。

\paragraph{校准、覆盖率与风险。}
预测区间覆盖率和区间宽度需要同时观察。过宽的区间容易达到覆盖率，却可能没有决策价值；过窄的区间看似精确，却会造成过度自信。对分类概率，可按置信度分箱比较预测概率和经验事件率；对回归区间，可按设备、工况和预测提前量分层检查。校准应在独立验证数据上完成，并在部署漂移后持续监测。

\paragraph{小样本故障率的后验决策。}
某设备在新工况下运行$20$次，观察到$2$次异常。若采用$\operatorname{Beta}(1,9)$先验，则后验为$\operatorname{Beta}(3,27)$，后验均值为$0.1$。直接使用样本比例得到的估计也是$0.1$，但二者对未来风险的表达不同：后验还包含小样本造成的宽区间。若把“安排人工复检”的成本记为$C_I$，把一次漏检损失记为$C_M$，则当后验异常概率满足
\begin{equation}
 p(\text{异常}\mid\mathcal D)C_M>C_I
\end{equation}
时才应触发复检。这样，概率输出通过成本函数转化为动作，而不是凭经验固定一个看似通用的阈值。

该例还说明先验不是隐藏的常数。若先验改为$\operatorname{Beta}(1,1)$，后验均值变为$3/22$；两种结果的差异提示当前数据不足以完全决定风险。工程报告应同时给出先验、后验和决策阈值，使不同团队可以复核结论。

\paragraph{回归预测区间的覆盖实验。}
假设模型对每台设备输出均值$\mu_i$和标准差$\sigma_i$，构造区间$[\mu_i-1.96\sigma_i,\mu_i+1.96\sigma_i]$。在验证集上统计真实值落入区间的比例，并分别按设备、环境温度和预测提前量分组。若整体覆盖率为$0.94$，但高温工况只有$0.71$，则总体数字掩盖了局部失效。可以在独立校准集上估计分位数修正因子，再在另一段时间上复验，避免把校准集上的偶然波动误认为改进。

区间宽度也应进入报告。两个模型可能具有相同覆盖率，但一个区间宽得无法指导维护，另一个区间更窄且只在少数工况失效。因而应联合绘制覆盖率、平均区间宽度和漏覆盖成本，并说明校准是否使用了未来数据。

\paragraph{后验预测检查的操作步骤。}
后验预测检查不是把训练数据重新拟合一次，而是从后验参数样本生成复制数据$\tilde y$，比较$\tilde y$与真实数据的统计量。例如可比较均值、方差、极值数量、故障持续时间和相邻残差相关性。若模型能匹配均值却无法产生真实的长尾，说明其似然假设不适合风险评估。工业案例中，检查统计量应优先选择与维护动作相关的量，而不是只选择容易展示的总体误差。

本章的后验、图模型和近似推断可分别参考下面的课程材料与官方概率编程教程，前者建立概念，后者展示可计算实现。
\section*{辅助学习材料}
\begingroup
\emergencystretch=3em
\begin{itemize}
\item \textbf{Probabilistic Graphical Models}；作者/平台：Stanford University CS228；英文；课程讲义与课程资源。重点学习 Bayesian networks、Markov networks、因子分解和推断主题，帮助理解本章的 HMM、MRF、消息传递与条件独立。\href{https://ermongroup.github.io/cs228/}{课程主页}。
\item \textbf{Bayesian Regression - Introduction (Part 1)}；作者/平台：Pyro 社区；英文；官方教程。阅读模型与先验、变分推断和后验预测部分，将点估计与参数不确定性作对照。\href{https://pyro.ai/examples/bayesian\_regression.html}{在线阅读}。
\item \textbf{【深度学习数学基础 23】统计推断与贝叶斯框架}；知乎；中文解说。先写先验、似然与后验，再复算本章一维更新例子；标题与主题据必应索引，原页受限。\href{https://zhuanlan.zhihu.com/p/1910026999415673453}{在线阅读}。
\item \textbf{贝叶斯推断概览}；CSDN；中文博客。比较参数后验与后验预测，注意噪声假设；标题与主题据必应索引，原页受限。\href{https://blog.csdn.net/qq\_41282102/article/details/84928605}{在线阅读}。
\item \textbf{基于GMM-HMM的语音处理全过程讲解-通俗易懂}；上传者：\_静静老师；bilibili；中文技术讲解。以语音为例对照本章隐马尔可夫模型，区分隐状态转移与观测分布，再比较工业状态估计的条件假设。2026年10月4日公开元数据播放4,404次；本次候选中选择主题直接相关的讲解，不宣称达到十万播放或全站最高；已核对标题与计数，未逐帧观看。\href{https://www.bilibili.com/video/BV1Rv4y1c7SA/}{观看视频}。
\end{itemize}
\endgroup
\paragraph{本章小结。}
贝叶斯更新表达参数在当前证据下仍有多少不确定性，概率图模型则明确状态、观测和标签之间保留哪些依赖。HMM 建模状态与观测的联合分布，MEMM 和 CRF 直接建模条件状态分布，MRF 将邻接关系推广到一般无向图。结构决定可以使用哪些推断算法，参数学习决定这些算法依据什么概率进行计算。即使推断在数学上精确，错误的观测假设或时间边界仍会使维护决策失效。后验预测检查、覆盖率与先验敏感性因此构成模型使用前的检验，而不是公式推导后的附加说明。

贯穿这些方法的共同步骤是先写联合或条件分布，再利用结构进行求和、积分或近似，最后在实际成本下解释预测。解析解提供清楚的参照，动态规划节省状态枚举，变分与采样处理无法解析的参数积分；各自的适用边界都来自前提，而非算法名称。小样本复检案例中的阈值还隐含复检能消除所计漏检损失的简化假设，若检查不完美或有后续成本，应扩充动作损失再取后验期望。

下一章继续使用概率，但不只给有限个系数设分布，还直接考虑可能的函数。线性高斯后验中用过的配方，将再次用于高斯过程的条件分布；核函数则说明两个输入有多相似、相应的函数值应有多大相关性。读者可以据此比较参数化与非参数化模型的灵活程度、先验假设和计算开销。
